\section*{Chapitre \ref{ch:g2:measuring-length} --- Mesurer les longueurs}

\begin{solution}{exo:g2:measuring-length:1}
Tout près~: se mettre dos à dos. Loin~: se \omterm{def:g2:measuring-length:measure}{mesurer} en \omterm{def:g2:measuring-length:units}{centimètres} et
envoyer le nombre --- un nombre peut voyager dans une lettre, une
comparaison dos à dos non.
\end{solution}

\begin{solution}{exo:g2:measuring-length:2}
Ni l'une ni l'autre ne s'est trompée. Les pas de Mara sont plus
courts~: il en tient donc davantage dans la même salle de classe. Des
\omterm{def:g2:measuring-length:measure}{unités} différentes donnent des nombres différents pour une même
longueur.
\end{solution}

\begin{solution}{exo:g2:measuring-length:3}
Le \omterm{def:g2:measuring-length:units}{centimètre} est le même pour tout le monde --- il n'appartient au
corps de personne. Une mesure en \omterm{def:g2:measuring-length:units}{centimètres} peut être vérifiée par
n'importe qui, avec n'importe quelle règle.
\end{solution}

\begin{solution}{exo:g2:measuring-length:4}
La fourmi~: \unit{cm} (et même moins~!)~; la porte~: \unit{m}~; la
cour de récréation~: \unit{m}~; la largeur de ce livre~: \unit{cm}.
\end{solution}

\begin{solution}{exo:g2:measuring-length:5}
Non --- il y a en général un petit espace vide entre le bord d'une
règle et son trait du zéro~: le crayon est donc un peu plus court que
$13$ \omterm{def:g2:measuring-length:units}{centimètres}. Tim a oublié de placer le \emph{trait du zéro}, et
non le bord, à l'extrémité du crayon.
\end{solution}

\begin{solution}{exo:g2:measuring-length:6}
Les réponses varient~; en général~: la main environ \qty{7}{cm} de
large, la cuillère environ \qty{15}{cm}, la tasse environ
\qty{9}{cm}. L'essentiel~: chaque nombre porte son \unit{cm}.
\end{solution}

\begin{solution}{exo:g2:measuring-length:7}
$100 + 100 = 200$~: la porte mesure \qty{200}{cm} de haut.
\end{solution}

\begin{solution}{exo:g2:measuring-length:8}
Les réponses varient~; la plupart des lits mesurent près de
\qty{2}{m}. Ce qui compte, c'est de comparer la réponse devinée au
nombre mesuré.
\end{solution}

\begin{solution}{exo:g2:measuring-length:9}
Le ruban de \qty{45}{cm} est le plus long~: $45 - 38 = 7$, donc de
\qty{7}{cm}. Sans nombres, les poser côte à côte trouve aussi le plus
long --- mais «~\qty{45}{cm}~» tient dans une lettre, et «~côte à
côte~» non.
\end{solution}

\begin{solution}{exo:g2:measuring-length:10}
Personne d'autre n'a les jambes de Papi~: un petit enfant qui
arpenterait le même jardin compterait peut-être $60$ petites
enjambées, et un visiteur ne pourrait pas vérifier «~$30$
enjambées~» sans Papi. Mesurée en \omterm{def:g2:measuring-length:units}{mètres} --- les mêmes pour tout le
monde --- la longueur du jardin peut être écrite, envoyée et vérifiée
par n'importe qui.
\end{solution}
